\section{Part D: unresolved gates and unchecked candidates}

\paragraph{Current scope after the original-boundary and local-orbit syntheses.}
Only general D.1, general D.2(ii), and their dependent residual L8
remain requested open rows. The checked local projective theorem
\ref{D:SNC-general-projective-quadratic-count} covers arbitrary SNC
exceptional graphs, including secondary residues, and
\ref{D:complete-local-projective-growth-all-dimensions} gives the
stated complete local dimensional and finite-normalization closure.
Theorem~\ref{D:local-divisorial-tame-projective-growth} also treats
every fixed complete normal local surface with an original reduced
divisorial tame boundary. Its constant is independent of rank,
coefficient field and horizontal root order; exceptional wild images
are retained.

Theorems~\ref{D:global-finite-puncture-projective-growth} and
\ref{D:global-original-boundary-projective-growth} prove the
unrestricted-period quadratic projective count for every fixed
normal proper surface with its original boundary convention.
The stated dimensional, controlled-centre and finite-normalization
closures are given by
\ref{D:global-projective-finite-puncture-closure} and
\ref{D:global-tame-projective-controlled-closure}.
Every defining-characteristic Lie simple-product quotient is
finitely generated in these scopes. Hereditary projective constants
linear in the degree of arbitrary finite covers are now proved for
every fixed bound on geometric Brauer period, as described below.
The unbounded-period hereditary projective count remains unproved.

The restricted alternating factorial count in
\ref{D:finite-puncture-local-orbit-factorial-count},
\ref{D:fixed-local-quotient-alternating-count}, and
\ref{D:alternating-dimensional-normalization-closure} uses the
whole fixed detecting-surface local groups. It requires the
square-sum of uncontrolled orbit sizes to be $O(n\log(n+1))$.
Orbits factoring through any chosen fixed finitely generated
local quotients may be omitted; in particular, nilpotent-image
and prime-to-$p$-image orbits may have arbitrary size.
Disconnected covers and the stated dimensional and nonnormal
closures are included. Unrestricted alternating multiplicities
remain open.

Theorems~\ref{D:alternating-compressed-local-orbits} and
\ref{D:original-boundary-matrix-alternating-count} enlarge that
criterion by prescribed small matrix models, including affine point
and projective partial-flag actions, and by original divisorial tame
boundaries on normal proper surfaces. The intrinsic enumeration in
\ref{D:intrinsic-small-linear-permutation-images} and
\ref{D:intrinsic-linear-orbit-alternating-count} also leaves uncharged
every whole local orbit whose image has faithful $\mathbb F_p$
dimension $D_N$ with $D_N^4=O(N\log(N+1))$. No structural choice
of an affine or flag action is required for this last criterion.
All these statements enlarge the same restricted alternating gate.
Corollary~\ref{D:primitive-defining-characteristic-orbit-models}
recognizes every stated primitive local action with nonalternating
defining-characteristic simple socle factors, together with a fixed
finite collection of other simple factors. It uses the known
primitive-order theorem and has faithful dimension
$O((1+\log N)^2)$. This recognition belongs to the same criterion.

Proposition~\ref{D:hereditary-wreath-permutation-count} and
Theorem~\ref{D:alternating-local-wreath-energy} now replace the
whole-orbit square by the minimum weighted tower energy
$\sum_{i\text{ unrestricted}}T_i m_i^2$ in the finite-puncture
scope. Modelled levels contribute only a factorial-order count.
Arbitrary-height towers of the stated affine and nonalternating
defining-characteristic primitive constituents therefore have zero
cost. The one-level unrestricted certificate recovers the earlier
square-sum criterion. This is one enlarged alternating criterion;
it does not extend the hereditary theorem to original-divisorial
local pairs or include unrestricted alternating socles.

Corollary~\ref{D:non-natural-alternating-primitive-models} further
recognizes the non-natural alternating-socle primitive actions.
The subset/product models have squared matrix rank at most $4N$;
the small-order branch has the checked polylogarithmic catalogue.
These are modelled levels of the same finite-puncture tower criterion.
Unbounded natural $A_N/S_N$ actions retain their unrestricted cost.

The Cartier recovery and generation statements
\ref{D:hereditary-Cartier-parameter-jets} and
\ref{D:hereditary-Cartier-generation-budget} are auxiliary.
Lemma~\ref{D:fixed-henselian-algebraization} and
Proposition~\ref{D:all-complete-cover-geometric-budgets} now remove
the former algebraization qualification from the geometric budgets:
one fixed henselian base works for all complete punctured covers.
These geometric inputs alone do not supply arbitrary-family
realization or the hereditary count. The subsequent ordinary
synthesis supplies the required object and comparison steps.

The ordinary-unit determinant list in
\ref{D:hereditary-ordinary-unit-determinants} and the actual finite-cycle
Hom presentations in
\ref{D:hereditary-ordinary-unit-finite-cycle-Hom} and
\ref{D:hereditary-generation-cycle-Hom-budget} have the stated
degree-linear budgets on both $bC$ and $bW$. They retain the full
nonreduced cycles and work over arbitrary coefficient algebras for
each fixed determinant choice. The determinant enumeration itself
concerns geometric ordinary unit systems, not twisted projective
determinants or arbitrary Artinian determinant variation.
Propositions~\ref{D:weak-cycle-Artinian-full-realization} and
\ref{D:hereditary-weak-cycle-relative-realization} remove the strong
cycle hypothesis from the conditional relative realization theorem.
Cohomological freeness, its stated Artinian base changes, and the
split second-neighbourhood boundary remain required. No degree-linear
presentation of these family conditions or hereditary count follows
from the finite complexes alone.
The exact intrinsic Hom-window results
\ref{D:intrinsic-Hom-fixed-loss-jets} and
\ref{D:hereditary-canonical-Hom-arrow-window} give prescribed
lower-cycle lifts for already factored canonical families: loss
three for ordinary Hom and $p+2$ for one Frobenius step, with the
stated threshold for the opposite bounded-pole arrow. Ordinary
isomorphisms lift globally. This does not impose the factorization
or realization hypotheses on arbitrary finite-cycle data, or make
independently lifted Frobenius relations compatible.

Lemma~\ref{D:strict-Cartier-separating-parameters} and
Proposition~\ref{D:uniform-comparable-reduction-cycle} choose one
strictly anti-nef separating parameter pencil on the fixed base.
On every specified relative-minimal cover resolution, the actual
Jacobian divisor gives a strong reduction cycle
$Z=\alpha C-\beta R_f$ with $C\le Z\le\alpha C$, for fixed
$\alpha,\beta$. Its energy, canonical intersection and fixed-level
scalar cohomology are linear in the cover degree. This aligns the
strong-cycle and parameter-jet constructions without an additional
cover-dependent depth. Arbitrary further blowups are not included.

The auxiliary refinement in
Lemma~\ref{D:linear-SNC-refinement-centres} and
Corollary~\ref{D:SNC-refinement-inherited-cycles} now gives an SNC
model with $O_R(t)$ centres, components, nodes and normalization
genus. The old actual cycles are pulled back: their complete
nonreduced cohomology, intersections, generation and specified
ordinary-family inverse poles are preserved. The controlled
non-SNC centres also give the stated Jacobian inequality.
This does not bound unsigned products of neighbouring cycle
multiplicities or by itself provide a projective presentation.
The bounded-period construction below uses ordinary full modules
on the original comparable-cycle model.

Six later auxiliary statements clarify the bounded-period route.
Lemma~\ref{D:full-bundle-cycle-degree} and
Proposition~\ref{D:canonical-order-Morita-conductor} give
degree-linear full-bundle budgets and actual order/comparison
maps with a fixed Cartier conductor. The scaled-separability
tensor and Newton formulas in
\ref{D:full-order-scaled-separability}--\ref{D:scaled-separable-action-Newton}
lift local algebra actions with an explicit loss of precision.
The geometric comparisons
\ref{D:stable-root-anchor-canonical-comparison}--\ref{D:bounded-period-generating-root-anchor}
give a bounded-rank generating anchor on a bipartite refinement.
Their lattice sandwiches do not prove generation or cohomology
bounds for arbitrary intermediate lattices. The additional
simultaneous Morita-unit and overlap effectivity is supplied by
Theorem~\ref{D:finite-Morita-Frobenius-object-effectivity};
no hereditary projective count follows from the six auxiliary
units alone.

Proposition~\ref{D:intrinsic-integral-unit-Hom-jets} gives exact
finite equations for integral formal unit maps between already
actual bounded-pole families. Local Frobenius contraction proves
effectivity and gluing; finite proper-cycle Hom functors give
finite presentation. For the canonical relative families satisfying
\ref{D:canonical-family-unit-Isom-openness}, every specified
Artinian unit isomorphism lifts uniquely and integrally, so unit
classes are open and closed. In the ordinary canonical case,
$c=p$ and $N=3$ suffice, with equations on $(p+3)C$.
These statements retain actual family existence, relative fullness
and the coflat/split-boundary hypotheses. They give no effective
coordinate or degree bound for a parameter locus containing all
objects. For ordinary objects, the following simultaneous object
construction and generated fixed-determinant comparison resolve
that remaining step without imposing coflat strata.

Theorem~\ref{D:arbitrary-cover-ordinary-hereditary-growth} now gives
\[
 \#\bigl(\operatorname{Hom}_{\rm cont}
 (\pi_1(U_A),\operatorname{GL}_r(\mathbb F_q))
       /\operatorname{GL}_r(\mathbb F_q)\bigr)
                         \le q^{A_Rtr^2}
\]
for every degree-$t$ connected finite punctured cover of a fixed
complete excellent normal local surface, with no tameness condition
on that cover. Proposition~\ref{D:finite-Frobenius-object-effectivity}
constructs actual compatible objects from the two finite inverse
equations, retaining the prescribed lower jet. The generated
fixed-determinant comparison in
\ref{D:weak-cycle-generated-determinant-lifting} and
\ref{D:generated-fixed-determinant-unit-openness} gives clopen
unit classes even for the noncanonical objects in this atlas.
Fixed determinant tuples, all-ring finite Hom complexes, and exact
boundary witnesses give $O_R(ftr^2)$ variables and equations of
bounded degree. The ordinary count is proved. Its supporting
statements are not separate formal-review movements; the
bounded-period projective extension uses the additional
Morita effectivity and determinant comparison below.

Lemma~\ref{D:arbitrary-field-character-scalar-H1} bounds the crossed
scalar characters by $q^{c_Rt}$, including wild field-cover order.
The exact graph/Jacobian filtration and finite-order deck action
give the bound without a local all-characteristics $H^1$ theorem.
The equivariant comparable model in
\ref{D:equivariant-cover-comparable-cycles} is the same quantitative
model, with no additional blowups. Theorem
\ref{D:arbitrary-cover-semilinear-hereditary-growth} consequently
proves $q^{A_R^{\rm sl}tr^2}$ conjugacy classes for every prescribed
field character. The field-cover degree $d$ cancels against the
short-cycle length $f/d$. No projective lifting is inferred.

Theorem~\ref{D:arbitrary-quotient-faithful-linear-crowns} now controls
the stated faithful-linear crowns on every complete normal surface
germ and for arbitrary permutation quotients. A normal subgroup of
bounded outer index must have the stated faithful linear model and
quadratic order comparison. The crown multiplicity is uniformly
polynomial in the socle order. The generator estimate retains
$\max\{2,d(P),c_R+2\}$, so no bound on an arbitrary quotient's
generator number has been inserted. This and the ordinary hereditary
count form one concrete enlargement of the controlled crown family;
the numerical and family inputs are not separate movements.

Corollary~\ref{D:arbitrary-quotient-faithful-semilinear-crowns}
also allows a faithful semilinear model for that bounded-index
normal subgroup. The fixed outer map prescribes its field
character, with no bound or tameness condition on the field-image
order. The same restriction and extension argument gives polynomial
socle multiplicity and retains the quotient term $d(P)$.
This is a consequence of the same arbitrary-cover hereditary
reduction, not an additional projective or full-generation claim.

In the earlier one-curve radially tame projective scope,
\ref{D:tame-projective-automorphism-crowns} and
\ref{D:normal-subgroup-projective-crowns} permit unbounded field
image; the classical and exceptional recognition results include
all defining-characteristic Lie families in that same scope.
Theorem~\ref{D:bounded-period-hereditary-projective-growth}
now proves $q^{A_{R,B}tr^2}$ ambient projective conjugacy classes
on every connected degree-$t$ punctured cover of a fixed complete
excellent normal local surface, for geometric Brauer period at
most a fixed $B$. The full scalar coefficient class is retained,
with no bound on $q-1$. Canonical full orders and Hom modules,
the actual Morita object construction, and the generated
fixed-determinant comparison in
\ref{D:generated-Morita-unit-Isom-openness} give an atlas with
$O_{R,B}(ftr^2)$ coordinates and bounded-degree equations.
The determinant lists are actual generated lists from canonical
inputs, and the retained finite marking preserves every original
punctured class. Relative root-stack hull flatness is not used.

Two further qualitative improvements do not complete the unbounded-period
atlas. Proposition~\ref{D:rank-period-generating-root-anchor} uses the
prime-to-characteristic local period-index theorem and one blowup of
each mixed node to obtain a generating twisted anchor of rank exactly
the geometric period $M$. Its coarse order has ordinary rank $M^2$;
rank-$r$ twisted bundles correspond to locally free order-modules of
rank $r/M$. This exact rank cancellation does not prove uniform
whole-cycle generation, syzygy or cohomology bounds for those modules.

Lemma~\ref{D:exceptional-infinitesimal-Picard-kernel} proves that a
specified Artinian punctured line comparison retains exactly its
closed-fibre exceptional divisor. Proposition
\ref{D:effective-formal-unit-family-class-openness} therefore removes
fixed determinant charts from qualitative unit-class openness for
already effective finite-type ordinary or fixed-order Morita families.
Its Serre twist is used only in comparison. It does not construct the
object family or give the required $O_R(ftr^2)$ coordinate budget.
These statements are auxiliary improvements, without another
hereditary count or an independent main-gate reduction.

Corollary~\ref{D:bounded-period-arbitrary-quotient-projective-crowns}
gives the corresponding faithful ordinary projective crown bound
over arbitrary permutation quotients, retaining the bounded
residual outer index, quadratic simple-order comparison and
$d(P)$. A sufficient multiplier condition is
$\gcd(r,q-1)\le B$; the stated compatible central-lift condition
is another. This is one bounded-period projective counting and
crown enlargement, not separate movements for each supporting
construction.

Lemma~\ref{D:arbitrary-field-character-H2} now also bounds the
actual twisted coefficient classes by $q^{c_Rt}$, including
field-cover order divisible by $p$. Its Cartan--Leray proof retains
the middle descent-cohomology term and the quotient correction in
the exceptional intersection lattice. The skew Morita construction
in \ref{D:skew-Morita-Frobenius-effectivity} retains the actual
gerbe and coherent deck functor. Its terminal action is the deck
pullback of the initial action, and its finite-frame comparison
uses the given deck cocycle. Theorem
\ref{D:bounded-period-semilinear-hereditary-projective-growth}
therefore proves the same $q^{A_{R,B}^{\rm ps}tr^2}$ bound for
projective representations with every prescribed field character,
provided the geometric Brauer period after restriction to the
field kernel is at most $B$. The short cycle has $f/a$ stages on
a cover of degree $ta$, so no extra factor $a$ remains in the
parameter bound. No order-$a$ matrix relation or division by $a$
is imposed. This is part of the same bounded-period enlargement.

Corollary~\ref{D:bounded-period-arbitrary-quotient-projective-semilinear-crowns}
extends that faithful crown consequence to arbitrary prescribed
field characters and nonnormal bounded-index model subgroups.
The actual scalar class is killed by $2$ whenever the specified
matrix model consists of semisimilitudes of a fixed nonzero
bilinear or quadratic form. This is an explicit algebraic model
condition, without an unchecked assertion about a named family.

Corollary~\ref{D:arbitrary-quotient-symplectic-orthogonal-crowns}
checks that condition for the exact natural symplectic and orthogonal
models, including the stated graph indices and the doubled coefficient
field in minus type. Corollary
\ref{D:global-arbitrary-quotient-symplectic-orthogonal-crowns}
gives the same arbitrary-quotient crown conclusion on a fixed normal
proper surface with finitely many omitted points. The full twisted
scalar class is killed by two before restriction to each actual local
field kernel. Both results retain $d(P)$ and are consequences of the
same bounded-period enlargement.

Example~\ref{D:unitary-unbounded-scalar-obstruction} explains why the
natural unitary form does not give an analogous uniform exponent:
explicit shift and clock matrices produce scalar classes of exact
unbounded odd order $m$ inside the natural projective unitary models.
The subsequent construction in
\ref{D:elliptic-cone-Heisenberg-covers} and
\ref{D:fixed-elliptic-cone-unitary-geometric-periods} realizes exact
unbounded geometric periods on the puncture of one fixed elliptic
cone. A valuation restriction identifies the actual descended cyclic
algebra and proves its exact period by the unramified norm criterion.
The resulting images are nilpotent Heisenberg groups. They are not
epimorphisms onto the ambient simple unitary groups and do not refute
the open hereditary count. This is a geometric obstruction to a
uniform-period shortcut, not a further main-gate reduction.

The numerical fixed-local curve lemma
\ref{D:curve-fixed-local-growth-numerical} has exponent proportional
to $1+g+|P|+\sum_x d_x$, where $d_x$ bounds the normalized
different of each full local projective type. Pulling back one
fixed detecting curve, Theorem
\ref{D:global-finite-puncture-hereditary-bounded-local-period}
therefore proves $q^{A_{S,T,B}tr^2}$ for ordinary projective classes
on arbitrary degree-$t$ covers of a fixed finite-puncture normal
proper surface, with bounded local geometric periods on the
fixed completed-source factors. Exact different transitivity,
wild Riemann--Hurwitz and summed local jet lengths keep every
curve cost linear in $t$. Each surface-cover factor is counted
once, even if it supplies several curve branches. This transfers
the same bounded-period enlargement to the stated global scope.

Theorem~\ref{D:global-finite-puncture-hereditary-semilinear-local-period}
now gives that global hereditary bound for every prescribed field
character as well. The affine-curve cyclic-band vanishing in
\ref{D:affine-curve-field-twisted-H2} gives a same-field semilinear
lift on the detector, and
\ref{D:local-field-character-crossed-H1} counts local scalar-lift
classes downstairs with exact cardinality $q_x-1$.
The numerical skew-curve lemma
\ref{D:curve-fixed-local-semilinear-growth-numerical} uses a
deck-linearized reference line pulled back from the base curve.
Genus and total jet depth on the field cover are $O(a)$ times
the original numerical budget, while only $f/a$ phases occur.
Its full local skew comparison preserves the actual field action.
The local geometric periods are measured after each local field
kernel; no ordinary Brauer class is assigned to a nonconstant
band on the base. This is the semilinear closure of the same
global bounded-period transfer.

Unbounded geometric period, projective crowns outside the proved
criteria, uncovered cross-characteristic families, unrestricted
alternating multiplicities and full finite generation, and
higher-dimensional boundaries outside the proved reductions
remain unresolved. The following record retains successive
approaches and their later closures; a historical intermediate
gate is not a new open item.

\subsection{D.1 and D.2(ii): a fixed singular compactification}
\begin{remark}[Unresolved in this run]\label{D:fixed-boundary-gate}
The remaining D.1 question concerns a normal proper nonprojective
$X$ in positive characteristic, and tameness tested only at the
codimension-one components of one fixed boundary. Outside the positive
cases of Proposition~\ref{D:fixed-boundary-positive}, we have not
proved finite generation. Proposition~\ref{D:boundary-neighborhood}
further reduces the gate to boundaries with no quasi-projective open
neighborhood. In particular, the one-point case of D.2(ii) is proved;
a finite omitted set in the residual case must exhibit a genuine
multiple-point obstruction to a common quasi-projective neighborhood.
This includes a regular $X$ with boundary
not having normal crossings. In D.2(ii), regular compactifications
are settled by purity; its residual normal case is singular.
Proposition~\ref{D:local-purity-boundary} further settles D.2(ii)
whenever every omitted local ring is regular of dimension at least
two or a complete intersection of dimension at least three.
In particular all codimension-at-least-three omitted subsets in a
local-complete-intersection neighborhood are settled, with arbitrary
singularities retained in the open.  This adds higher-dimensional
cases without asserting purity at a singular surface point.
In dimension two the residual field must have positive
transcendence degree over $\mathbb F_p$, by
Corollary~\ref{D:finite-field-surfaces}.
For normal surfaces, Corollary~\ref{D:finite-cone-fixed-boundary}
settles D.1 whenever each non-$\mathbb Q$-factorial boundary
singularity is isolated in the boundary and belongs to one of
the proved local classes: log canonical, formally positively
graded, a single smooth exceptional curve satisfying the finite
cone-cover criterion, a full formal retraction, or either explicit
genus-two germ.  The numerical condition $-C^2>2g(C)-2$ suffices
in genus at least two; at equality it suffices that
$\omega_C\otimes N_{C/\widetilde S}$ is torsion, in particular
that the germ is $\mathbb Q$-Gorenstein with any index.
Thus a residual D.1 surface must have a non-$\mathbb Q$-factorial
boundary singularity lying on a boundary curve, or an isolated
point outside all these classes.  Retained singularities are
unrestricted.
Proposition~\ref{D:finite-local-groups} also settles D.2(ii) for
strongly F-regular compactifications and, more generally, whenever
the punctured strict-local groups along the omitted set have
uniformly bounded finite order. In the residual finite-point case,
at least one such local group must therefore be infinite.
Proposition~\ref{D:finite-local-images-arbitrary} sharpens this
necessary condition for arbitrary omitted sets, without a uniform
order bound: at least one punctured strict-local group at a closed
omitted point must have infinite image in the global $\pi_1(U)$.
Corollary~\ref{D:klt-threefold-boundary} supplies a further
higher-dimensional class: proper normal threefolds in
characteristic $p>5$ with klt singularities near the arbitrary
codimension-two omitted set.  The retained singularities are
unrestricted.  A residual threefold in this characteristic
therefore has an omitted non-klt point.  The group is nevertheless
small, and the maximal pronilpotent quotient
of every open subgroup is finitely generated, for an arbitrary
omitted set of codimension at least two, by
Proposition~\ref{D:codim-two-small}. Proposition~\ref{D:fixed-boundary-small}
proves both restrictions for every fixed-boundary tame group in
D.1 as well. Thus an abelian or pro-$p$ quotient cannot refute
finite generation, and every fixed finite group supports only
finitely many torsors in the relevant category over every finite
cover. The remaining issue is a uniform generator bound across
all finite quotients, not finiteness for each chosen finite target.
Example~\ref{D:pronilpotent-obstruction} explains why this
restriction does not decide the full profinite group.
Proposition~\ref{D:uniform-H1-obstruction} strengthens the
obstruction: the same non-finitely-generated group has
$\dim H^1(G,M)\leq\dim M$ for every finite module and
Schreier-style cohomology bounds for every open subgroup.
Proposition~\ref{D:cohomology-crown-criterion} identifies the
precise two uniform estimates that would suffice: a linear
rank bound for $H^1$ and a polynomial bound for nonabelian
crown multiplicities.  The all-field linear estimate is now proved in dimension two by
Theorems~\ref{D:punctured-surface-all-characteristics-H1} and
\ref{D:tame-surface-all-characteristics-H1}, with the stated
Schreier bounds for every open subgroup. Hence the maximal
prosolvable quotients are finitely generated. The remaining surface
criterion concerns nonabelian crown multiplicities.
Theorem~\ref{D:global-quadratic-representation-bound} bounds
defining-characteristic representation classes by $q^{Cr^2}$ for
proper-surface punctures, including the stated nonnormal extension;
Theorem~\ref{D:uniform-local-representation-growth} and
Corollary~\ref{D:uniform-local-growth-all-dimensions} give the local
analogue in every dimension. Theorem~\ref{D:tame-parabolic-fixed-rank-growth}
gives $q^{C_r}$ for the original-divisor tame surface quotient;
Theorem~\ref{D:tame-global-quadratic-representation-bound} strengthens
this to $q^{Cr^2}$ with a constant independent of rank and coefficient field.
Example~\ref{D:quadratic-growth-alternating-obstruction} shows that
even an all-field quadratic bound combined with linear cohomology
does not imply finite generation abstractly. The exceptional-genus
bound of Corollary~\ref{D:surface-residual-genus-bound} is an
additional geometric constraint, not a bound for all crown
multiplicities. Unrestricted-dimensional nonprojective global
reduction also remains open. The
two-singularity mutation surfaces of
Proposition~\ref{D:mutation-surfaces} are also settled, even when
their boundary has no common quasi-projective neighborhood.
Corollary~\ref{D:schroer-elliptic} additionally settles the
three-point surface with trivial Picard group when the curve in
Schr\"oer's construction is elliptic.  Its proof uses
Proposition~\ref{D:lc-elliptic-tame}, not finiteness of the
punctured local groups: every quasi-\'etale cover is tame above a
positive-genus exceptional curve of a log canonical surface.
Proposition~\ref{D:single-curve-residue-tameness} removes the
log canonical and $\mathbb Q$-Gorenstein assumptions for the
single smooth elliptic curve case.  Theorem~\ref{D:surface-mixed-boundary}
and Corollary~\ref{D:lc-surface-boundary} prove the stronger mixed
result: every finite omitted set consisting of locally
$\mathbb Q$-factorial points, single smooth elliptic contractions,
or arbitrary log canonical points has finitely generated punctured
fundamental group.  Singularities retained in the open are
unrestricted, and the stated finite-normalization extension
treats nonnormal surfaces as well.
The proof uses Lemma~\ref{D:lc-exceptional-configurations}:
a minimal exceptional fibre at a log canonical point is a rational
tree, one elliptic curve, or a rational nodal cycle.
Proposition~\ref{D:rational-tree-Qfactorial} makes the tree case
locally $\mathbb Q$-factorial in positive characteristic, without
assuming a rational or log terminal singularity.  We retain those
points on a projective partial resolution and use tameness only
at the elliptic and cyclic boundary divisors.  Every remaining
normal surface counterexample must therefore omit a point which
is non-$\mathbb Q$-factorial, non-log-canonical, and outside the
single smooth elliptic class.  Theorem~\ref{D:formal-cone-mixed-boundary}
and Corollary~\ref{D:weighted-graded-boundary}
further exclude every finite omitted set consisting of
$\mathbb Q$-factorial, log-canonical, or formally positively graded
points.  Theorem~\ref{D:finite-cone-mixed-boundary} further adds
the finite cone-cover, formal-retraction, and explicit genus-two
classes.  A residual single smooth exceptional curve must have
genus at least two, must admit no full formal retraction, and
must satisfy $-C^2\leq2g(C)-2$.  At equality,
$\omega_C\otimes N_{C/\widetilde S}$ must be nontorsion, so the
germ is not $\mathbb Q$-Gorenstein.  Moreover
Proposition~\ref{D:finite-jet-cone-cover} excludes every retraction
of $C_r$ for which $r\geq2$ and $r(-C^2)>2g(C)-2$.
Thus if $-C^2>g(C)-1$, the first-neighborhood class must be
nonzero.  The finite-cover obstruction criterion of
Lemma~\ref{D:finite-formal-cone-cover} must fail as well.  No Cartier-index gate remains for the global log canonical
finite-generation statement.
Theorem~\ref{D:surface-map-to-curve} gives a different global
reduction: every finite puncture of a normal proper surface
admitting a nonconstant morphism to a proper curve is settled,
with no local restriction on the omitted singularities and
without requiring a smooth generic fibre.  A residual finite
surface puncture must therefore have no such morphism.
Propositions~\ref{D:vertical-fibration-exact} and
\ref{D:vertical-fibration-multiplicity} also settle opens with
vertical boundary meeting every fibre on an appropriate proper
birational model.  Their full base-inertia images are finite
by the fixed fibre multiplicities; infinite source inertia is
not confused with its image in the required quotient.
More generally,
Proposition~\ref{D:finite-exceptional-wild} proves that every
projective modification of a remaining counterexample must have
a boundary divisor with infinite wild image in the original
group.  Infinite source local groups or finite exceptional wild
inertia alone are insufficient obstructions.
The theorem for the valuation-tame fundamental group of every finite-type
scheme over an algebraically closed field, in
\cite[section 2.6]{AHLS2026}, tests boundary divisors in all normal
compactifications. Its hypothesis is stronger and cannot be substituted.
Likewise, replacing $X$ by a projective modification can create
exceptional boundary divisors: covers of $X\setminus Y$ may be wildly
ramified there, even when the original $Y$ has codimension at least two.
This is the exact failed step in a proposed reduction by Chow's lemma.
Example~\ref{D:fixed-vs-valuation} exhibits the extra wild inertia
explicitly on a quotient of $\mathbb P^1\times\mathbb P^1$ in
characteristic two. Its full group is finite, so it refutes the
identification of the two tameness conditions, not finite generation.
The punctured-local-spectrum finite-generation conjecture recorded in
\cite[Conjecture 1.1]{Suzuki2023} indicates a related difficulty;
Theorem 1.2 there proves a nilpotent-quotient statement, not the full
fundamental-group assertion. We assert no equivalence between that
local conjecture and the global questions here.
\par
The previously residual explicit three-point surface with
$\operatorname{Pic}(S)=0$ in
\cite[sections 3.1--3.2, author's PDF p.~5]{Schroer1999D}
is now settled in every positive genus by
Corollary~\ref{D:schroer-all-genera}.  The two single-curve
contractions have split formal neighborhoods
$\mathcal O_C(-c_1)$, so Proposition~\ref{D:graded-cone-tame}
proves radial tameness, including separability of the residue
extension.  The projective partial contraction of $A\cup R_2$
then leaves only those two tame divisorial boundary components.
No condition is imposed at its remaining codimension-two point.
Thus this construction is no longer a residual test or a
candidate counterexample.  The distinct horizontal reduction
of Proposition~\ref{D:Schroer-horizontal-tame} remains valid.

For more general positive-genus exceptional curves, the
weighted inequality of
Proposition~\ref{D:single-curve-residue-tameness} still does not
control fierce residue extensions without a split formal-cone
structure.  Example~\ref{D:fierce-small-excess} rules out a
numerical shortcut: in each fixed characteristic, Galois surface
function-field extensions can have inseparable residue and
$(\delta-e+1)/e=2/p^r\to0$.  These are local extensions, not
quasi-\'etale covers of a contracted proper surface.  The
remaining geometric finite-generation problem therefore needs
a further global restriction, or an actual compatible geometric
realization with unbounded finite-quotient generator numbers.
Example~\ref{D:wild-genus2-single-curve} now supplies a genuine
quasi-\'etale fierce cover of an affine contracted surface:
one smooth genus-two exceptional curve with self-intersection
$-1$ supports a quadratic cover with $e=1$, inseparable residue
degree two, and different exponent one.  It refutes universal
single-curve tameness by an actual geometric example.
It does not refute finite generation, and no unbounded family
of finite quotient generator numbers is obtained from it.
Its source resolution has one genus-two curve of square $-2$,
and Corollary~\ref{D:wild-genus2-pro2} proves that the strict-local
pro-two quotients of the source and target are respectively
trivial and $\mathbb Z/2\mathbb Z$.  Thus the displayed wild
quadratic step cannot simply be iterated.  That calculation
alone does not control nonabelian wild images.  The separate
finite formal cone construction now does:
Corollary~\ref{D:wild-genus2-finite-inertia} bounds the full
exceptional wild images by four at the source and eight at
the target, with no covering-group restriction.  These explicit
germs are consequently removed from the remaining local tests
by Theorem~\ref{D:finite-cone-mixed-boundary}.  Smallness and
pronilpotent finite generation alone would not supply those
bounds, as Example~\ref{D:pronilpotent-obstruction} proves.
Theorem~\ref{D:surface-map-all-boundaries} now excludes every
normal proper surface admitting a nonconstant morphism to a proper
curve from both remaining questions, with arbitrary original
horizontal and vertical boundary in D.1.  A rational pencil on
a resolution does not suffice: exceptional curves over omitted
points can be horizontal, and the original condition imposes no
tameness at those new horizontal valuations.
In arbitrary dimension, Corollaries~\ref{D:vertical-boundary-all-dimensions}
and~\ref{D:projective-generic-normalization} further exclude vertical
boundaries over a proper curve and arbitrary boundaries whose
geometric generic normalization is projective.  For vertical
codimension-at-least-two boundary the conclusion concerns the full
open fundamental group.  Proposition~\ref{D:generic-fibre-generation}
records the exact more general geometric generic-group hypothesis.
Proposition~\ref{D:projective-base-generation} extends that reduction
to normal projective bases of arbitrary dimension, bounding the
inseparable residue degree as well as the wild ramification index.
Corollary~\ref{D:projective-image-corollaries} in particular handles
vertical boundaries over any positive-dimensional projective image,
and all boundaries when the normalization of the reduced geometric
generic fibre is projective, including generic fibre dimension at
most one.  Theorem~\ref{D:iterated-fibration-generation}
propagates this conclusion, including all finite normal covers,
along any tower of proper morphisms with projective reduced
geometric generic normalizations; fibre dimensions at most one
suffice at every stage.
Propositions~\ref{D:log-regular-boundary}
and~\ref{D:generically-finite-projective-target} also exclude
proper log regular pairs with their interior boundary, and
normal proper schemes admitting a generically finite morphism
toward a projective target.  These are recognition criteria;
a log structure without log regularity or an alteration in the
opposite direction does not supply their hypotheses.
\end{remark}



% Audited post-checkpoint13 bounded-central-order reduction
\paragraph{The remaining projective and crown gate.}
Theorem~\ref{D:surface-tame-cross-characteristic-H2} proves
rank-linear second cohomology in coefficient characteristics distinct
from the ground characteristic. Corollary~\ref{D:surface-tame-central-obstructions}
controls the central lifting classes. For every fixed bound on the natural
centre, Corollary~\ref{D:surface-bounded-centre-simple-product} removes
all ranks of the stated type-A simple-product family. The same conclusion
holds for nonnormal proper-surface punctures by
Corollary~\ref{D:nonnormal-bounded-centre-projective-growth}.
The constant is allowed to depend on that centre bound. The later
Brauer-zero theorem below removes this dependence for the stated global
surface projective representations with zero geometric Brauer class.
Unbounded central order with nonzero Brauer class in that setting,
alternating multiplicities, uncovered cross-characteristic Lie families
whose orders are divisible by the ground characteristic, semilinear actions, general nonabelian
crowns, and global reduction outside the finite-puncture and controlled-centre
classes below remain unresolved.


% Local and natural classical recognition after the same criterion
\paragraph{Local and named-family scope.}
Corollary~\ref{D:strict-local-open-cohomology-finite} gives the
invertible-coefficient finiteness input in arbitrary local dimension.
Corollary~\ref{D:local-bounded-scalar-simple-product} applies the
bounded-scalar criterion to complete equicharacteristic normal punctures
in every dimension. Corollaries~\ref{D:bounded-centre-unitary-simple-product}
and \ref{D:symplectic-orthogonal-simple-product} verify the named natural
classical families, including their scalar kernels and small exceptions.
This recognition does not turn a simple-product statement into a bound
for general crowns and does not cover the remaining cross-characteristic
families or unbounded linear and unitary centres.


% Auxiliary exceptional core and Frobenius saturation
\paragraph{What the new exceptional estimates leave open.}
Proposition~\ref{D:exceptional-arithmetic-genus-bound} bounds the
arithmetic genus of the full reduced exceptional fibre and its graph cycles
linearly in cover degree. Proposition~\ref{D:exceptional-degree-linear-core}
also bounds the core components, their negative squares, genera and all
attachments to the remaining rational $(-2)$-blocks. These earlier
estimates alone leave the vertex count unrestricted; the later
Proposition~\ref{D:punctured-cover-linear-resolution} bounds it for
covers unramified outside an isolated point. Lemma
\ref{D:normalization-Frobenius-saturation} supplies a degree-independent
pole bound for the inverse Frobenius map on the actual finite
normalization lattice. These auxiliary bounds do not yet supply a
hereditary object atlas with parameter count linear in cover degree,
a factorial cover-count bound, or full nonabelian-crown control.


% Post-checkpoint14 finite omitted-locus dimensional reduction
\paragraph{The finite omitted-locus reduction and its limit.}
Theorem~\ref{D:generic-surface-finite-punctures} constructs one fixed
normal proper surface over an algebraically closed extension whose
punctured fundamental group surjects onto the original group whenever
the omitted locus is finite. Corollaries
\ref{D:normal-finite-puncture-quantitative-transfer} and
\ref{D:nonnormal-finite-puncture-quantitative-transfer} transfer the
linear first-cohomology and quadratic representation bounds, the
prosolvable conclusion and the bounded-scalar simple-product conclusions
in every dimension, including the stated nonnormal case. Positive-dimensional
omitted loci still need a global reduction, and full nonabelian-crown
control remains open even on punctured surfaces. No second-cohomology
injection is inferred from the group surjection.

\paragraph{Why the auxiliary repairs do not close that gate.}
Example~\ref{D:superperfect-quadratic-obstruction} adds rank-linear
second cohomology to the abstract non-generation obstruction.
Proposition~\ref{D:contract-zero-cost-blocks} and
Lemma~\ref{D:RDP-uniform-local-generators} compress the long rational
$(-2)$-blocks geometrically and bound full local puncture generators.
Example~\ref{D:RDP-branch-gluing-obstruction} shows that branch
deletion needs additional control; the inherited different and global
extension constraints remain relevant and are not refuted by that example.
Example~\ref{D:Chow-surface-boundary-obstruction} prevents an unrestricted
ample-slicing shortcut when the original omitted locus has positive
dimension. None of these auxiliary facts supplies the missing hereditary
cover-count atlas or full finite generation.


% Post-checkpoint14 controlled exceptional centres
\paragraph{An explicit extension to positive-dimensional boundaries.}
Theorem~\ref{D:generic-surface-controlled-boundary} gives the surface
reduction if a fixed normal projective modification has only point
images for exceptional divisors centred inside the omitted locus.
Its original-divisorial-boundary quotient maps surjectively to the
original quotient, using generic transversality and absolute Abhyankar.
For codimension-at-least-two boundary the surface has finite punctures.
Corollary~\ref{D:nonnormal-controlled-boundary-quantitative-transfer}
descends the quantitative conclusions through finite normalization.
This narrows the dimensional gate further, but the existence of such
a modification is not asserted for arbitrary pairs. Positive-dimensional
exceptional centres lying in the boundary remain uncontrolled.


% Post-checkpoint14 full generation through codimension-one projective models
\paragraph{A further full-generation class.}
Proposition~\ref{D:codimension-one-projective-model} proves full finite
generation when a dense open of a normal proper variety is identified
with an open of a projective normal model whose complement has
codimension at least two. No codimension condition is required on
the original side. Corollary~\ref{D:complete-toric-fixed-boundary}
applies this to every complete toric variety in arbitrary characteristic
and every closed boundary, including noninvariant boundaries. The stated
nonnormal descent is also full finite generation. A general proper
variety is not asserted to admit such a model; the surface problem and
the remaining unrestricted-dimensional cases therefore remain open.
Proposition~\ref{D:weighted-exceptional-chain-compression} further
compresses exceptional graphs. The later
Proposition~\ref{D:punctured-cover-linear-resolution} also bounds the
constant-multiplicity chains for covers unramified outside an isolated
surface point. Neither result supplies a cover-count atlas.


% Post-review09 finite normal closure and punctured graph bound
\paragraph{Two further limits of the present conclusions.}
Proposition~\ref{D:finite-normal-tame-descent} closes the proved
full-generation classes under finite dominant normal images, including
inseparable finite maps. Its image-index estimate does not apply to
arbitrary merely generically finite modifications with new exceptional
boundary divisors. Proposition~\ref{D:punctured-cover-linear-resolution}
bounds the full exceptional graph size linearly in cover degree for
isolated punctures, using the Jacobian cost of a constant-multiplicity
chain. The remaining atlas problem includes formal-neighborhood and
gluing parameters; a graph-size bound alone does not settle it.
These are auxiliary or closure results, not further formal-review
movement after review09.


% Post-checkpoint15 canonical and weighted-genus bounds
\paragraph{Further auxiliary control, with the atlas still missing.}
Proposition~\ref{D:cover-uniform-canonical-defect} gives one fixed
base ideal annihilating the canonical defect and the resolution's first
cohomology for every finite punctured cover of a complete equicharacteristic
normal surface. Its geometric genus is linear in cover degree.
Corollary~\ref{D:cover-thickening-linear-cohomology} bounds the full
cohomology of each fixed pulled-back exceptional thickening linearly.
Proposition~\ref{D:cover-uniform-Noether-different} improves the
source-factorization ideal to its first power, including the stated
Artinian coefficient families. Lemma~\ref{D:exceptional-weighted-genus}
bounds the sum of squared arithmetic genera divided by negative
self-intersection. These conclusions do not control the number of global
cover objects, a hereditary atlas with the required degree bound, or the
remaining nonabelian crowns. They do not count as new formal-review movement.


% Post-checkpoint15 local Eisenstein and RDP degree bounds
\paragraph{Local object data and the remaining global comparison.}
Proposition~\ref{D:local-Eisenstein-linear-atlas} presents all separable
extensions of a Laurent series field with fixed degree and different
exponent using a linear number of coefficients. It includes an explicit
finite witness for pointed isomorphisms over arbitrary coefficient rings.
Corollary~\ref{D:local-branch-data-linear-complexity} controls the number
and degree of charts when total different is linear in cover degree.
Proposition~\ref{D:RDP-full-sheaf-degree} bounds the total degree of a
reflexive module's full pullback on an entire rational-double-point
block by its rank. Neither theorem bounds global cover objects or supplies
the compatibility with adjacent exceptional curves needed for a uniform
hereditary atlas. These are auxiliary results, not new formal-review movement.


% Post-checkpoint15 fixed-fibre, energy and anchored-order results
\paragraph{The sharpened local and gluing statements.}
Corollary~\ref{D:unpointed-Eisenstein-witness} supplies the same finite
isomorphism witness for continuous unpointed formal-disc maps, including
nilpotent translations. Proposition~\ref{D:fixed-formal-curve-fibre-bound}
bounds each fixed-full-formal-data fibre on a genus-$h$ curve by
$(n!)^{2h}n^{2D}$. Its common-lift argument does not impose a high-degree
or Galois hypothesis. The missing step is an effective bound on the
additional locus of local data arising from the fixed surface problem;
fibre cardinality does not bound that locus.
Proposition~\ref{D:exceptional-canonical-energy} and
Corollary~\ref{D:punctured-cover-linear-energy} give a degree-linear
canonical-energy budget and control genus against effective intersection
numbers after elimination of rational blocks. They do not present the
cover objects. Example~\ref{D:elliptic-cone-order-generators} refutes
constant module-generator bounds for arbitrary Azumaya anchors on one
fixed punctured elliptic cone: the exact number is $3r$. A better anchor
is not excluded. None of these auxiliary conclusions closes a requested
row or proves the remaining nonabelian-crown bound.


% Post-checkpoint16 integral cycle budget
\paragraph{An integral cycle without denominator loss.}
Lemma~\ref{D:integral-antinef-energy-cycle} constructs an integral
anti-nef exceptional cycle by global lattice minimization and zero-cost
repairs on the zero-canonical components. Its energy is linear in cover
degree, and it has strictly positive conormal degree on the nonrational
core. Corollary~\ref{D:punctured-cover-antinef-cycle} bounds cohomology on
every fixed multiple of the full cycle, including its nilpotent layers.
These statements remove a numerical denominator loss. They do not prove
formal determination, boundary-framed compatibility, or the required
object-presentation bound, and are not new main-gate movement.


% Post-checkpoint16 Brauer-zero main-gate reduction
\paragraph{Unbounded scalar order with zero geometric Brauer class.}
Theorem~\ref{D:Brauer-zero-projective-representation-growth} proves
$q^{Cr^2}$ growth for all defining-characteristic
$\operatorname{PGL}_r(\mathbf F_q)$ representations of a fixed normal
proper surface original-boundary quotient whose geometric Brauer class
vanishes. Their Kummer lifting class need not vanish, and no bound on
the scalar-centre order is assumed. A fixed finite collection of line
bundles represents every prime-to-$p$ Picard quotient; scalar digits
replace the growing root by one-step Frobenius tensor cycles.
Lemmas and propositions~\ref{D:fixed-tensor-boundary-comparison} and
\ref{D:finite-tensor-cyclic-quadratic-atlas} supply the uniform boundary,
Euler, cohomology and Artinian quadratic-presentation estimates.
This is one strict reduction of the unbounded-centre projective gate:
the bounded-geometric-period closure below further restricts this
part of the global normal surface problem to unbounded geometric period. It does not prove a
simple-product or general-crown conclusion for the residual classes.
The remaining cross-characteristic families and alternating
multiplicities require their separate arguments. The local surface
Brauer-zero extension is stated below.
Corollaries~\ref{D:Brauer-zero-detecting-surface-transfer} and
\ref{D:Brauer-zero-normalization-transfer} now give the detecting-surface
and finite-normalization closure of this same reduction. They require
zero Brauer class only after each fixed surface pullback, with no
injection in degree-two cohomology asserted. General D.1,
D.2(ii) and residual L8 remain unresolved.


% Post-checkpoint16 RDP Hom repairs
\paragraph{What boundary framing does and does not cost.}
Lemma~\ref{D:RDP-full-Hom-vanishing} proves unframed Hom vanishing
for full sheaves on rational double points and a finite-subscheme
bound by rank product times length. The exact standard two-end
bound on an $A$-chain is at most $2rs$, independently of its length,
by Proposition~\ref{D:A-chain-two-end-Hom-bound}. However
Example~\ref{D:A-chain-internal-framing-obstruction} has ranks one,
one transverse internal branch and framed dimension $m$ on an
$A_{3m-1}$ chain, even for fixed-characteristic Frobenius-periodic
reflexive modules. Thus attachment count alone does not suffice.
These individual-module statements do not yet provide an effective
global family atlas or the needed nonabelian count.


% Post-checkpoint16 local Brauer-zero and finite-jet closure
\paragraph{The complete local Brauer-zero extension.}
Theorem~\ref{D:local-Brauer-zero-projective-growth} gives the same
$q^{Cr^2}$ count on the puncture of every complete excellent normal
equicharacteristic surface ring with algebraically closed coefficient
field, without a finite-type algebraization hypothesis.
Lemma~\ref{D:local-Picard-prime-to-p-generators} supplies fixed
generators modulo every prime-to-$p$ integer.
Proposition~\ref{D:local-fixed-tensor-cyclic-atlas} lifts the entire
tensor source from fixed jets using source identity factorizations;
the source itself need not be MCM or coefficient-flat. This is the
local scope of the same Brauer-zero reduction. Its all-dimension fixed-surface closure is
Corollary~\ref{D:local-Brauer-zero-surface-transfer}. It leaves
nonzero geometric Brauer classes after the fixed surface pullbacks,
local settings outside the stated hypotheses, and full generation unresolved.

\paragraph{Finite jets and their actual-bundle families.}
Proposition~\ref{D:A-chain-finite-jet-framing} replaces whole-branch
compatibility on an $A$-chain by fixed finite jets when attachments
are pairwise distinct and transverse. The inverse-Cartan reflection
identity and Schur complement give total length $O(n)$ in
Corollary~\ref{D:A-block-total-jet-budget}.
Proposition~\ref{D:RDP-Artinian-full-bundle-framing} proves rigidity
and the exact jet obstruction over Artinian coefficients for actual
vector bundles on the fixed resolution with full closed fibres.
It does not identify arbitrary relative MCM pullbacks with such
bundles, vary the resolution, control coincident-branch contacts,
or supply an effective hereditary object atlas. These three units
are auxiliary and do not create another movement event.

\paragraph{The selected simple-product consequence.}
Corollary~\ref{D:Brauer-zero-selected-simple-product} proves finite
generation of the product defined by the selected Brauer-zero simple
quotient kernels, whenever the fixed faithful projective embeddings
satisfy $|H|\ge q^{\epsilon r^2}$. It includes no unselected kernel
of the same simple types and gives no general crown bound. This is
a consequence of the single Brauer-zero count, not separate movement.


% Post-checkpoint17 sparse DVR and local transfer closure
\paragraph{The local dimensional closure and sparse extension charts.}
Corollary~\ref{D:local-Brauer-zero-surface-transfer} gives the
surface-detected Brauer-zero count on complete excellent normal
equicharacteristic punctures in every dimension. This is a closure
of the earlier count, not a new independent reduction.
Proposition~\ref{D:sparse-DVR-linear-atlas} gives two-generator
finite flat charts and exact finite unpointed Isom witnesses over
arbitrary coefficient rings, including imperfect residue fields.
The normalized different condition gives linear coordinate counts
and polynomial equation degrees in
Corollary~\ref{D:sparse-DVR-normalized-different-complexity}.
The residue polynomial is fixed. Its global poles and those of the
other coefficients remain uncontrolled, as does the degree of the
constrained surface-extension locus. No full cover count follows.


% Post-checkpoint17 bounded geometric Brauer period
\paragraph{Every fixed geometric period bound.}
Theorem~\ref{D:bounded-geometric-Brauer-period} proves the same
$q^{C_Br^2}$ count when the geometric Brauer class has order at most
a fixed integer $B$, on the stated normal proper surfaces and
complete local surfaces. The finite set of realized classes is killed
on one fixed normal cover obtained from witness representations.
Restriction to that cover and a coset-representative fibre bound
reduce the assertion to the Brauer-zero theorem. Scalar-centre order
is still unrestricted. The constant and the chosen cover may depend
on $B$; no bound uniform in $B$ is proved.
Corollary~\ref{D:bounded-Brauer-period-selected-simple-product}
gives the corresponding selected-kernel product consequence.
This is a finite-cover closure of the same single main reduction,
not another movement event. The projective route still requires
unbounded geometric Brauer period, and the remaining alternating,
cross-characteristic and general-crown gates are unchanged.


% Post-checkpoint17 generation charts and one-curve residues
\paragraph{Generation and a fixed-determinant object chart.}
Theorem~\ref{D:reflexive-pullback-uniform-generation} proves
uniform generation after one exceptional-cycle twist, using
automatic adjunction on every effective subcycle, including
nilpotent layers and the point-ideal Tor term.
Proposition~\ref{D:hereditary-fixed-determinant-cycle-chart}
combines the fixed Noether-different ideal with the energy
estimate and the cyclic Frobenius determinant recurrence.
It gives an actual $O(nr^2)$-coordinate section chart for a fixed
determinant on one finite exceptional cycle. Its energy-family
and complete-base hypotheses are explicit. The unresolved step
is finite-thickening realization and determination of the unit
objects, all Frobenius arrows and specified isomorphisms over
Artinian bases, with the same complexity. The chart alone
gives no hereditary representation count.

\paragraph{One smooth exceptional curve and a period-degree cover.}
Proposition~\ref{D:single-curve-Brauer-residue-cover} identifies
the prime-to-$p$ Brauer torsion with curve-character cohomology
and kills a period-$m$ class by a cyclic degree-$m$ cover
etale over the entire resolution. Genus and normal degree
scale as stated in that proposition. Constants on these
varying covers still need uniform control. A splitting cover
alone does not supply a global Clifford rank drop or a
formal weight-block decomposition. These three statements
are auxiliary; no main gate or requested status changes.


% Post-checkpoint18 exact canonical module jets
\paragraph{Actual module jets from two proper thickenings.}
Lemma~\ref{D:canonical-pullback-uniform-jet-loss} and
Proposition~\ref{D:canonical-module-two-level-recovery}
recover the fixed jets of every actual reflexive module from the
image of sections on two consecutive parameter-ideal thickenings.
The fixed level is independent of cover degree and module rank,
and both section and cohomology dimensions are $O(nr)$.
These thickenings can have embedded zero-dimensional structure.
No base-change property of the section images or realization
of an arbitrary chart family is inferred. Effective object and
arrow diagrams with the required equation-degree budget remain
the missing counting input. These two results are auxiliary.


% Post-checkpoint18 geometric realization inputs
\paragraph{Known reduction cycles and the qualified relative realization.}
Proposition~\ref{D:energy-controlled-reduction-cycle} shows that
the established degree-linear energy cycles already satisfy the
Kahn--Gavran hypotheses. The exact known pointwise correspondence
is recorded in Theorem~\ref{D:reduction-cycle-full-module-correspondence},
with the actual first-step marking action and characteristic-free
proof scope checked. It gives no prescribed-arrow lifting.
Theorem~\ref{D:Artinian-full-neighbourhood-realization} supplies
actual Artinian module families when an explicit second-neighbourhood
bundle, cohomological freeness and base change, and a split injective
boundary are retained. Corollary~\ref{D:Artinian-canonical-two-level-jets}
recovers their module jets and proves base change of that recovered
quotient, without asserting base change of the raw section image.
Effective equations for the required strata and compatible arrows,
with the desired variable and degree bounds, remain unresolved.

\paragraph{Exact one-curve depth and the residual twisted category.}
Proposition~\ref{D:single-curve-Cartier-parameter-ideal} chooses
one separating parameter pair inside the fixed defect ideal whose
pullback is a fixed multiple of the sole exceptional curve.
Corollary~\ref{D:single-curve-Cartier-module-jet-recovery} therefore
identifies the actual recovery thickenings on every etale-resolution
cover with fixed Cartier multiples, independently of cover degree.
Propositions~\ref{D:residual-twisted-spectral-category} and
\ref{D:residual-twisted-Frobenius-transport} describe the residual
twisted category at composite scalar order and its precise
Frobenius transport, retaining the possibly nonidentity final
torsor automorphism. Neither these curve statements nor the
Cartier-depth bridge provides the missing global normal-direction
Frobenius atlas. No unbounded-period count or main-gate reduction
is inferred from these eight inputs.


\paragraph{Finite map models and twisted normal depth.}
Proposition~\ref{D:fixed-determinant-finite-Hom} constructs actual
polynomial models for maps, composition, equality and inverse witnesses
on a fixed Cohen--Macaulay curve, for all coefficient algebras. Its
constant-degree formulas and hereditary coordinate budget do not give
bounded equations for the required coflat object strata or arbitrary
prescribed surface arrows.
Lemma~\ref{D:twisted-connected-anchor-minimal-slopes} constructs a
mixed-band formal anchor and normalized twisted lattices when the
residual spectral torsor is connected. Its Euler quotient has precisely
the spectral rank, giving a slope window of width independent of the
period. Proposition~\ref{D:twisted-period-independent-normal-depth}
then gives fixed original-curve pole and map-lifting depths, including
actual Artinian families, and quadratic Hom cohomology. A finite
graded object/arrow atlas with uniformly bounded equation degrees is
still missing. These three inputs leave every strongest main gate
unchanged; no unbounded-period representation count follows from them.


\paragraph{Post-checkpoint19 quantitative inputs and their scope.}
Lemma~\ref{D:weighted-affine-component-bound} and its two
normal-layer corollaries bound components of actual weighted
polynomial schemes; their small background block permits
polynomial coefficient degrees without adding a logarithm
to the entire quadratic variable budget.
Lemmas and propositions
\ref{D:weighted-fixed-residual-bundle-deformations}--\ref{D:weighted-fixed-determinant-normal-presentation}
construct actual fixed-determinant normal bundle and map
families. Lemma~\ref{D:compressed-minor-background-charts}
and Proposition~\ref{D:moving-line-polynomial-background}
supply rational moving-line and cohomology backgrounds on
explicit exact-rank chart schemes. They distinguish geometric
coverage of those strata from arbitrary Artinian lifting.
Proposition~\ref{D:one-curve-coefficient-root-gerbe}
extends the actual coefficient gerbe at its geometric period,
including its Kummer part. Lemma
\ref{D:meromorphic-unit-jet-contraction} and Proposition
\ref{D:finite-jet-unit-Isom} recover actual unit isomorphisms
from paired map/inverse jets and prove clopen class fibres
under their stated effectivity hypotheses. These ten units
alone do not prove an unbounded-period representation count:
that further conclusion requires controlled normalized anchor
coefficients and a complete audited synthesis. They do not
settle any general finite-generation row or the residual crown
and simple-family gates.


\paragraph{The one-curve unbounded-period count.}
Lemma~\ref{D:controlled-one-curve-gerbe-anchors} now supplies
actual normalized coefficient data with polynomial normal-layer
complexity, including the residual cyclic coherence and the final
spectral translation. Theorem~\ref{D:one-curve-unbounded-period-projective-count}
combines it with the effective weighted atlas and the actual unit-Isom
construction to bound all projective rank-$r$ representations by
$q^{K_Rr^2}$. This applies to complete excellent normal local surfaces
with one smooth connected exceptional curve, without a fixed bound on
the geometric Brauer period. Formal review13 records this as one
strict reduction of the projective gate. Other exceptional
configurations, the residual global projective problem, alternating
multiplicities, cross-characteristic families, general nonabelian
crowns, and full finite generation remain unresolved.


\paragraph{A restricted crown gate is now closed.}
Theorems~\ref{D:etale-resolution-hereditary-linear-count},
\ref{D:unramified-permutation-linear-crowns}, and
\ref{D:unramified-field-permutation-crowns} bound the
multiplicity and generator number of the specified
permutation and field-permutation crowns when their
quotient extends etale over the one-curve resolution.
The faithful linear/order inequality is part of the
hypotheses. The field-action order need not be bounded:
Lemma~\ref{D:unramified-semilinear-determinant-budget}
and Theorem~\ref{D:etale-resolution-semilinear-count}
keep its determinant and phase costs within the
index-linear budget. This is one restricted crown
reduction. Quotients ramified above the exceptional
curve, other outer actions and the remaining simple
families are not covered; no full-generation conclusion
for the punctured fundamental group follows.


The restricted crown reduction also permits radial tame ramification:
Theorem~\ref{D:tame-resolution-hereditary-linear-count} gives
$q^{C_Rtr^2}$ linear counts on every radially tame degree-$t$ cover,
and Theorem~\ref{D:tame-resolution-semilinear-count} gives the
corresponding bound for any radially tame fixed field character.
Corollary~\ref{D:tame-field-permutation-crowns} consequently controls
the stated faithful field-permutation crowns whenever their quotient
is tame along the exceptional curve. This enlarges the same single
post-review13 crown reduction. Arbitrary wild permutation covers,
other outer actions, the residual simple families and general
D.1/D.2(ii) remain open. The radial quotient used here is an auxiliary
quotient for this resolution; the original D.1 boundary is unchanged.


\paragraph{Ordinary projective heredity and bounded-field automorphism crowns.}
Theorem~\ref{D:tame-hereditary-projective-count} supplies
$q^{C_Rtr^2}$ ordinary projective counts on all degree-$t$ radially
tame covers of the one-curve resolution, using actual fixed partial
resolutions. Lemma~\ref{D:radial-tame-cyclic-H2} bounds the full
twisted coefficient obstruction on the tame quotient without an
invalid degree-two inflation through the normal wild subgroup.
Proposition~\ref{D:nonsplit-automorphism-wreath-lifts} gives the exact
nonabelian lifting bijection without assuming that the outer
extension splits. Theorem~\ref{D:bounded-field-automorphism-crowns}
then controls the stated projective automorphism crowns when the
graph quotient and the stabilizer field-image order are bounded.
The field size, projective rank and diagonal part remain unrestricted.
These statements belong to the one review14 crown reduction.

In the earlier reduction, the unbounded-field projective count in
\ref{D:bounded-graph-projective-crown-reduction} was an explicit
hypothesis, discharged by \ref{D:tame-projective-semilinear-count}. The geometric step that had remained was to extend the actual
coefficient gerbe coherently with the deck and band actions on the
field cover and to retain the index-linear weighted complexity.
Choosing an arbitrary root-gerbe extension at the geometric period
does not prove this compatibility. No such descent is inferred from
the ordinary projective count or the obstruction-class bound.

\paragraph{Exceptional trees: geometry and map depth.}
Lemmas~\ref{D:SNC-tree-Brauer-geometry} and
\ref{D:SNC-tree-residual-category} give exact residues, actual
coefficient-gerbe extensions, component spectral categories and
reduced node gluing for a simple-normal-crossings exceptional tree.
Lemmas~\ref{D:SNC-tree-minimal-marked-windows} and
\ref{D:SNC-tree-uniform-Hom-depth} give uniform marked slope windows,
original-cycle pole bounds and normal Hom depth. They are auxiliary:
they do not by themselves yield an effective object atlas, its unit
class comparison, or an unbounded-period representation count for
this larger exceptional configuration.


\paragraph{The exceptional-tree projective gate is now closed.}
Lemmas and propositions
\ref{D:SNC-tree-small-marked-source}--\ref{D:SNC-tree-finite-unit-comparison}
now supply the missing actual component atlas, thick-node pasting and
finite unit-class comparison. Theorem
\ref{D:SNC-tree-projective-quadratic-count} proves
$q^{C_Rr^2}$ for every fixed complete excellent normal local surface
with a regular projective SNC exceptional-tree resolution, without
bounding its geometric Brauer period. This is one strict extension
beyond the one-smooth-curve case after review14. It does not give
hereditary-index constants for trees. At that stage arbitrary graphs and the unbounded-field comparison
were still unresolved; both are now proved in their stated scopes.
At that stage the residual global projective problem was also open;
the later finite-puncture synthesis below closes precisely that scope.
No full finite-generation conclusion is inferred.


\paragraph{Coherent unbounded-field descent is now proved in the one-curve scope.}
Lemmas and propositions
\ref{D:coherent-full-root-coefficient-gerbe}--\ref{D:compressed-skew-Frobenius-unit-jets}
extend the actual twisted-band coefficient gerbe with its full deck
coherence, choose one residual sector without losing punctured objects,
and compress the extra root-character layers. Theorem
\ref{D:tame-projective-semilinear-count} consequently proves the
index-linear projective semilinear bound with unrestricted field-image
order. Corollary~\ref{D:tame-projective-automorphism-crowns} removes
that boundedness condition from the stated projective crown family.
This is one further strict reduction after review14. The cover and
field character must still be radially tame over the one-curve
resolution, and the faithful projective embedding, order comparison
and bounded graph hypotheses remain. Wild outer quotient covers,
other exceptional configurations and general finite generation are
not settled by these proofs.


\paragraph{Arbitrary SNC graphs and all complete local punctures.}
The two geometry statements
\ref{D:SNC-general-Brauer-geometry}--\ref{D:SNC-general-residual-category}
retain graph-cycle residues and the noncommuting node algebra.
The actual weighted presentation and unit comparison in
\ref{D:SNC-general-weighted-presentation} yield
\ref{D:SNC-general-projective-quadratic-count}. The new all-dimensional
closure \ref{D:complete-local-projective-growth-all-dimensions}
uses a fixed slice, projective SNC resolution and finite-proper torsor
descent, and does not assume Brauer pullback injectivity.
These results enlarge one local-projective reduction. They do not
alone give the later global count, which also uses the fixed detecting
curve and its fixed-local-type count. An index-linear atlas for
arbitrary finite local covers remains unproved.


\paragraph{Named Lie families and the tame hypothesis.}
Corollary~\ref{D:exceptional-defining-characteristic-crowns}
completes the defining-characteristic Lie-family recognition,
including Suzuki and Ree. This is part of the same one-curve
radially tame crown reduction. Example
\ref{D:one-curve-wild-swap-cover} gives an alternative construction,
for every genus at least two, of the failure already exhibited in
\ref{D:wild-genus2-single-curve}. Corollary
\ref{D:wild-swap-exact-image} gives its exact wild-image order two
and its finite-generation bound. These do not remove the radial
tameness hypothesis from the general crown theorem.

\paragraph{Global projective counts with the original boundary.}
The different estimate in
\ref{D:different-Frobenius-invariant-lattice} controls each fixed
local type uniformly in its splitting degree. The fixed detecting
curve and three curve statements
\ref{D:curve-integral-unit-jet-comparison}--
\ref{D:curve-fixed-local-projective-growth} give the integral unit
comparison and quadratic count, with prescribed Artinian Isom lifts.
The local-pair cohomology, horizontal residue and marked-root atlas
then prove \ref{D:local-divisorial-tame-projective-growth}, with no
horizontal-root-order loss. The fixed original local sources in
\ref{D:tame-detecting-curve-local-pairs} give
\ref{D:global-original-boundary-projective-growth} for every original
surface boundary, without a global Brauer-vanishing or bounded-period
assumption. The dimensional, normalization and defining-characteristic
simple-product consequences belong to this one reduction.
Arbitrary-cover hereditary constants and general crowns remain open.

\paragraph{Restricted alternating multiplicities.}
The local permutation argument pays for the ambiguity of permutation
bases: a modular permutation representation need not determine its
underlying finite set. The exact resulting bound in
\ref{D:local-permutation-orbit-energy}, together with the fixed-formal
curve fibre theorem and the detecting curve, proves the factorial
bound in \ref{D:finite-puncture-local-orbit-factorial-count}.
Corollary \ref{D:fixed-local-quotient-alternating-count} removes all
orbits whose action factors through one of finitely many fixed
finitely generated local quotients. The maximal pronilpotent and
pro-prime-to-$p$ quotients are eligible known examples; their orbit
sizes are unrestricted. Corollary
\ref{D:alternating-dimensional-normalization-closure} includes
disconnected covers and the stated dimensional and normalization
closures. These form one restricted alternating-gate reduction.
The remaining whole-local orbit square-sum is still unbounded.
Curve-inertia
orbits alone are insufficient for this argument: their subspaces
need not extend as subrepresentations of the whole completed-local
group whose representation count is used.
